\chapter{Przegląd technologii chmurowych i narzędzi pomiarowych}
\label{ch:przeglad_technologii}

Rozdział niniejszy przedstawia szczegółową specyfikację oraz analizę technologiczną komponentów wchodzących w skład produkcyjnego środowiska badawczego. Wszystkie wybrane rozwiązania zostały dobrane pod kątem ich natywnego wsparcia dla architektur mikroserwisowych oraz możliwości precyzyjnej rejestracji metryk wydajnościowych w warunkach dynamicznych zmian obciążenia sieciowego.

\section{Platforma orkiestracji Google Kubernetes Engine (GKE)}
\label{sec:gke_architektura}
Podstawowym środowiskiem uruchomieniowym dla badanych mikroserwisów chmurowych typu SaaS jest platforma \textit{Google Kubernetes Engine (GKE)} działająca w trybie \textit{Standard} w ramach chmury publicznej \textit{Google Cloud Platform (GCP)}. Odstąpiono od wstępnych założeń implementacji lokalnej w oparciu o dystrybucje \textit{Minikube} lub \textit{k3s}, ponieważ lokalne środowiska zwirtualizowane nie odwzorowują w sposób dostateczny rzeczywistych opóźnień sieciowych warstwy routingu chmurowego, fizycznego czasu alokacji zasobów dyskowych oraz dynamicznej charakterystyki działania płaszczyzny kontrolnej (ang. \textit{Control Plane}).

Architektura zarządzana GKE zapewnia pełną izolację komponentów sterujących od puli węzłów wykonawczych (ang. \textit{Worker Nodes}). Płaszczyzna kontrolna realizuje zadania związane z utrzymaniem stanu klastra poprzez rozproszony rejestr \texttt{etcd}, harmonogramowaniem zasobów (\texttt{kube-scheduler}) oraz orkiestracją replik za pomocą mechanizmu \texttt{kube-controller-manager}. 

Węzły wykonawcze klastra zostały oparte na instancjach maszyn wirtualnych \textit{Google Compute Engine} typu \texttt{e2-standard-4}. Każda z maszyn fizycznych dysponuje 4 wirtualnymi rdzeniami procesora (vCPU) oraz 16~GB pamięci RAM, co gwarantuje stabilność metryczną i eliminuje interferencję zasobową (ang. \textit{noisy neighbor effect}). Środowiskiem uruchomieniowym kontenerów jest \texttt{containerd} działający pod kontrolą systemu operacyjnego \textit{Container-Optimized OS (COS)}. Szczegółowe zestawienie parametrów środowiska GKE przedstawiono w tabeli \ref{tab:gke-specyfikacja}.

\begin{table}[ht]
\centering
\caption{Specyfikacja techniczna produkcyjnego środowiska badawczego GKE}
\label{tab:gke-specyfikacja}
\begin{tabular}{|l|l|}
\hline
\textbf{Komponent infrastruktury} & \textbf{Konfiguracja i parametry techniczne} \\ \hline
Typ maszyn wirtualnych węzłów    & \texttt{e2-standard-4} (Google Compute Engine) \\ \hline
Zasoby pojedynczego węzła        & 4 vCPU (architektura x86\_64), 16 GB RAM \\ \hline
Wersja płaszczyzny kontrolnej    & Kubernetes v1.29.x (Regular Release Channel) \\ \hline
Środowisko wykonawcze kontenerów & \texttt{containerd} v1.7.x \\ \hline
System operacyjny węzłów         & \textit{Container-Optimized OS} zarządzany przez Google \\ \hline
Topologia sieciowa klastra       & Sieć \textit{VPC-native} oparta na aliasach adresów IP \\ \hline
\end{tabular}
\end{table}

\section{Automatyczny system skalowania poziomego (HPA)}
\label{sec:hpa_mechanizm}
Mechanizm \textit{Horizontal Pod Autoscaler (HPA)} odpowiada za dynamiczne dostosowywanie liczby replik kontenerów w ramach obiektów typu \texttt{Deployment} na podstawie analizy sygnałów telemetrycznych. Algorytm decyzyjny HPA wykonuje cykliczne obliczenia (domyślnie co 15 sekund) w oparciu o ogólny wzór matematyczny:

\begin{equation}
R_{desired} = \lceil R_{current} \times \frac{M_{current}}{M_{desired}} \rceil
\label{eq:hpa_wzor}
\end{equation}

gdzie:
\begin{itemize}
    \item[$R_{desired}$] -- wyznaczona, pożądana liczba replik,
    \item[$R_{current}$] -- aktualna liczba uruchomionych Podów,
    \item[$M_{current}$] -- bieżąca średnia wartość rejestrowanej metryki wydajnościowej,
    \item[$M_{desired}$] -- zdefiniowana przez administratora wartość progowa metryki targetowej.
\end{itemize}

W specyfikacji standardu API \texttt{autoscaling/v2} mechanizm HPA wspiera skalowanie przy użyciu trzech odmiennych typów metryk:
\begin{enumerate}
    \item \textbf{Metryki zasobowe (Resource Metrics):} wyznaczane jako procentowe zużycie CPU lub pamięci RAM w odniesieniu do wartości deklaratywnej \texttt{resources.requests}. Sygnały te są dostarczane przez komponent \texttt{metrics-server}.
    \item \textbf{Metryki niestandardowe (Custom Metrics):} parametry aplikacyjne powiązane bezpośrednio z obiektami Kubernetes (np. intensywność ruchu sieciowego RPS), pobierane z zewnętrznych systemów monitorowania poprzez rozszerzenie \texttt{Custom Metrics API}.
    \hline
    \item \textbf{Metryki zewnętrzne (External Metrics):} parametry niezwiązane bezpośrednio z obiektami klastra, pobierane z zewnętrznych usług chmurowych (np. długość kolejki wiadomości Pub/Sub).
\end{enumerate}

\section{Zintegrowany stos monitorujący kube-prometheus-stack}
\label{sec:telemetria_stos}
Akwizycja danych pomiarowych w środowisku chmurowym realizowana jest w sposób w pełni zautomatyzowany przy użyciu zintegrowanego stosu \textit{kube-prometheus-stack}, zainstalowanego w klastrze GKE za pomocą menedżera pakietów \textit{Helm}. System rezygnuje z nieefektywnej i podatnej na błędy ręcznej agregacji danych pomiarowych z logów terminala.

\subsection{Baza danych Prometheus i mechanizmy SD}
\label{subsec:prometheus}
Sercem systemu telemetrii jest baza danych szeregów czasowych \textit{Prometheus (TSDB)}. Kolekcjonowanie danych odbywa się za pomocą modelu ciągnionego (\textit{pull-model}), w którym serwer monitorujący wykonuje asynchroniczne zapytania HTTP do punktów końcowych \texttt{/metrics} z interwałem próbkowania wynoszącym $\Delta t = 5\text{ s}$. 

Automatyczne wykrywanie celów pomiarowych (ang. \textit{Service Discovery}) realizowane jest za pomocą mechanizmu \texttt{kubernetes\_sd\_configs}, który analizuje adnotacje kontenerów i dynamicznie rejestruje instancje mikroserwisów. Infrastruktura zbiera pełne metryki systemowe z demonów \textit{cAdvisor} (statystyki zużycia CPU i pamięci kontenerów na poziomie jądra Linux) oraz \textit{kube-state-metrics} (statusy obiektów orkiestracji).

\subsection{Prometheus Adapter v0.12.0}
\label{subsec:adapter}
Komponent \textit{Prometheus Adapter} pełni rolę mostu translacyjnego pomiędzy bazą danych Prometheus a natywnym mechanizmem HPA klastra GKE. Rejestruje się on w architekturze jako usługa \texttt{APIService}, rozszerzając standardowe API Kubernetes o ścieżkę \texttt{/apis/custom.metrics.k8s.io/}. Adapter pobiera chwilowe wartości liczników za pomocą języka zapytań \textit{PromQL}, przekształca je (np. stosując funkcję \texttt{rate()} w oknie czasowym) i serwuje w formacie zrozumiałym dla algorytmu autoskalera.

\subsection{Platforma wizualizacyjna Grafana}
\label{subsec:grafana}
Platforma \textit{Grafana} służy jako centralny punkt prezentacji danych empirycznych oraz weryfikacji przebiegu eksperymentów. Usługa Grafana została wystawiona na zewnątrz klastra GKE przy użyciu publicznego mechanizmu równoważenia obciążenia chmury GCP (\texttt{type: LoadBalancer}), co umożliwia bezpieczny dostęp do dedykowanego panelu \textit{SaaS Wydajność Overview} i automatyczny eksport danych telemetrycznych.

\section{Narzędzie testów obciążeniowych Grafana k6}
\label{sec:k6_charakterystyka}
Do generowania deterministycznego obciążenia sieciowego wykorzystano otwartoźródłowe narzędzie \textit{Grafana k6} w wersji \texttt{v0.48+}. Narzędzie to zostało napisane w języku \textit{Go}, natomiast scenariusze testowe definiowane są w języku \textit{JavaScript (ES6)}. Przewagą systemu \textit{k6} nad klasycznymi rozwiązaniami (np. Apache JMeter) jest wysoka wydajność jednowątkowa oraz natywne wsparcie dla asynchronicznych modeli wykonywania operacji wejścia-wyjścia, co pozwala na generowanie tysięcy współbieżnych użytkowników (\textit{Virtual Users -- VU}) bez wywoływania zjawiska saturacji zasobów maszyny pomiarowej.